% Options for packages loaded elsewhere
\PassOptionsToPackage{unicode}{hyperref}
\PassOptionsToPackage{hyphens}{url}
\PassOptionsToPackage{dvipsnames,svgnames,x11names}{xcolor}
\documentclass[
  11pt,
  a4paper,
]{article}
\usepackage{xcolor}
\usepackage[margin=1in]{geometry}
\usepackage{amsmath,amssymb}
\setcounter{secnumdepth}{-\maxdimen} % remove section numbering
\usepackage{iftex}
\ifPDFTeX
  \usepackage[T1]{fontenc}
  \usepackage[utf8]{inputenc}
  \usepackage{textcomp} % provide euro and other symbols
\else % if luatex or xetex
  \usepackage{unicode-math} % this also loads fontspec
  \defaultfontfeatures{Scale=MatchLowercase}
  \defaultfontfeatures[\rmfamily]{Ligatures=TeX,Scale=1}
\fi
\usepackage{lmodern}
\ifPDFTeX\else
  % xetex/luatex font selection
  \ifLuaTeX
    \usepackage{luaotfload}
    \directlua{luaotfload.add_fallback("mainfontfallback",{
      "NotoSerifBengali:mode=harf"
    })}
  \fi
  \setmainfont[,RawFeature={fallback=mainfontfallback}]{Latin Modern Roman}
  \ifLuaTeX
    \usepackage{luaotfload}
    \directlua{luaotfload.add_fallback("sansfontfallback",{
      "NotoSansBengali:mode=harf"
    })}
  \fi
  \setsansfont[,RawFeature={fallback=sansfontfallback}]{Latin Modern Sans}
  \ifLuaTeX
    \usepackage{luaotfload}
    \directlua{luaotfload.add_fallback("monofontfallback",{
      "NotoSansBengali:mode=harf"
    })}
  \fi
  \setmonofont[,RawFeature={fallback=monofontfallback}]{DejaVu Sans Mono}
\fi
% Use upquote if available, for straight quotes in verbatim environments
\IfFileExists{upquote.sty}{\usepackage{upquote}}{}
\IfFileExists{microtype.sty}{% use microtype if available
  \usepackage[]{microtype}
  \UseMicrotypeSet[protrusion]{basicmath} % disable protrusion for tt fonts
}{}
\makeatletter
\@ifundefined{KOMAClassName}{% if non-KOMA class
  \IfFileExists{parskip.sty}{%
    \usepackage{parskip}
  }{% else
    \setlength{\parindent}{0pt}
    \setlength{\parskip}{6pt plus 2pt minus 1pt}}
}{% if KOMA class
  \KOMAoptions{parskip=half}}
\makeatother
\usepackage{color}
\usepackage{fancyvrb}
\newcommand{\VerbBar}{|}
\newcommand{\VERB}{\Verb[commandchars=\\\{\}]}
\DefineVerbatimEnvironment{Highlighting}{Verbatim}{commandchars=\\\{\}}
% Add ',fontsize=\small' for more characters per line
\usepackage{framed}
\definecolor{shadecolor}{RGB}{248,248,248}
\newenvironment{Shaded}{\begin{snugshade}}{\end{snugshade}}
\newcommand{\AlertTok}[1]{\textcolor[rgb]{0.94,0.16,0.16}{#1}}
\newcommand{\AnnotationTok}[1]{\textcolor[rgb]{0.56,0.35,0.01}{\textbf{\textit{#1}}}}
\newcommand{\AttributeTok}[1]{\textcolor[rgb]{0.13,0.29,0.53}{#1}}
\newcommand{\BaseNTok}[1]{\textcolor[rgb]{0.00,0.00,0.81}{#1}}
\newcommand{\BuiltInTok}[1]{#1}
\newcommand{\CharTok}[1]{\textcolor[rgb]{0.31,0.60,0.02}{#1}}
\newcommand{\CommentTok}[1]{\textcolor[rgb]{0.56,0.35,0.01}{\textit{#1}}}
\newcommand{\CommentVarTok}[1]{\textcolor[rgb]{0.56,0.35,0.01}{\textbf{\textit{#1}}}}
\newcommand{\ConstantTok}[1]{\textcolor[rgb]{0.56,0.35,0.01}{#1}}
\newcommand{\ControlFlowTok}[1]{\textcolor[rgb]{0.13,0.29,0.53}{\textbf{#1}}}
\newcommand{\DataTypeTok}[1]{\textcolor[rgb]{0.13,0.29,0.53}{#1}}
\newcommand{\DecValTok}[1]{\textcolor[rgb]{0.00,0.00,0.81}{#1}}
\newcommand{\DocumentationTok}[1]{\textcolor[rgb]{0.56,0.35,0.01}{\textbf{\textit{#1}}}}
\newcommand{\ErrorTok}[1]{\textcolor[rgb]{0.64,0.00,0.00}{\textbf{#1}}}
\newcommand{\ExtensionTok}[1]{#1}
\newcommand{\FloatTok}[1]{\textcolor[rgb]{0.00,0.00,0.81}{#1}}
\newcommand{\FunctionTok}[1]{\textcolor[rgb]{0.13,0.29,0.53}{\textbf{#1}}}
\newcommand{\ImportTok}[1]{#1}
\newcommand{\InformationTok}[1]{\textcolor[rgb]{0.56,0.35,0.01}{\textbf{\textit{#1}}}}
\newcommand{\KeywordTok}[1]{\textcolor[rgb]{0.13,0.29,0.53}{\textbf{#1}}}
\newcommand{\NormalTok}[1]{#1}
\newcommand{\OperatorTok}[1]{\textcolor[rgb]{0.81,0.36,0.00}{\textbf{#1}}}
\newcommand{\OtherTok}[1]{\textcolor[rgb]{0.56,0.35,0.01}{#1}}
\newcommand{\PreprocessorTok}[1]{\textcolor[rgb]{0.56,0.35,0.01}{\textit{#1}}}
\newcommand{\RegionMarkerTok}[1]{#1}
\newcommand{\SpecialCharTok}[1]{\textcolor[rgb]{0.81,0.36,0.00}{\textbf{#1}}}
\newcommand{\SpecialStringTok}[1]{\textcolor[rgb]{0.31,0.60,0.02}{#1}}
\newcommand{\StringTok}[1]{\textcolor[rgb]{0.31,0.60,0.02}{#1}}
\newcommand{\VariableTok}[1]{\textcolor[rgb]{0.00,0.00,0.00}{#1}}
\newcommand{\VerbatimStringTok}[1]{\textcolor[rgb]{0.31,0.60,0.02}{#1}}
\newcommand{\WarningTok}[1]{\textcolor[rgb]{0.56,0.35,0.01}{\textbf{\textit{#1}}}}
\usepackage{longtable,booktabs,array}
\usepackage{calc} % for calculating minipage widths
% Correct order of tables after \paragraph or \subparagraph
\usepackage{etoolbox}
\makeatletter
\patchcmd\longtable{\par}{\if@noskipsec\mbox{}\fi\par}{}{}
\makeatother
% Allow footnotes in longtable head/foot
\IfFileExists{footnotehyper.sty}{\usepackage{footnotehyper}}{\usepackage{footnote}}
\makesavenoteenv{longtable}
\setlength{\emergencystretch}{3em} % prevent overfull lines
\providecommand{\tightlist}{%
  \setlength{\itemsep}{0pt}\setlength{\parskip}{0pt}}
% paper_preamble.tex — shared LaTeX preamble for the PCMA/SI paper set.
%
% Compiled with lualatex (NOT pdflatex) because several papers contain
% Bengali script: the CLL paper realises one cognitive state into
% English, Bengali and Banglish, so the Bengali output is EVIDENCE and
% cannot be transliterated away without destroying the point being made.
%
% Font selection is NOT done here. pandoc >= 3 has native
% `mainfontfallback` / `monofontfallback` variables that emit the
% luaotfload fallback registration in the correct position relative to
% unicode-math; doing it by hand in this header lands after unicode-math
% and breaks font loading outright. See tools/build_papers.sh.

\usepackage{amsmath}
\usepackage{amssymb}
\usepackage{mathtools}   % \xleftrightarrow etc. used in cll.md §5
\usepackage{booktabs}

% ---------------------------------------------------------------------
% Unicode symbols appearing in PROSE rather than inside $...$.
%
% The papers were written as markdown and read as markdown, so they use
% real characters — "the ω-boost", "α-child", "ΔR coherence" — instead of
% LaTeX macros. Under unicode-math those codepoints are math-mode-only
% and abort the build ("\symbballowed only in math mode"). Mapping each
% to \ensuremath keeps the markdown readable on GitHub AND compilable,
% which is the whole point of keeping one source.
% ---------------------------------------------------------------------
\usepackage{newunicodechar}

% Greek, lower case
\newunicodechar{α}{\ensuremath{\alpha}}
\newunicodechar{β}{\ensuremath{\beta}}
\newunicodechar{γ}{\ensuremath{\gamma}}
\newunicodechar{δ}{\ensuremath{\delta}}
\newunicodechar{ε}{\ensuremath{\varepsilon}}
\newunicodechar{η}{\ensuremath{\eta}}
\newunicodechar{θ}{\ensuremath{\theta}}
\newunicodechar{κ}{\ensuremath{\kappa}}
\newunicodechar{λ}{\ensuremath{\lambda}}
\newunicodechar{μ}{\ensuremath{\mu}}
\newunicodechar{π}{\ensuremath{\pi}}
\newunicodechar{ρ}{\ensuremath{\rho}}
\newunicodechar{σ}{\ensuremath{\sigma}}
\newunicodechar{τ}{\ensuremath{\tau}}
\newunicodechar{φ}{\ensuremath{\varphi}}
\newunicodechar{χ}{\ensuremath{\chi}}
\newunicodechar{ψ}{\ensuremath{\psi}}
\newunicodechar{ω}{\ensuremath{\omega}}
\newunicodechar{ν}{\ensuremath{\nu}}
\newunicodechar{ξ}{\ensuremath{\xi}}
\newunicodechar{ζ}{\ensuremath{\zeta}}
\newunicodechar{ι}{\ensuremath{\iota}}

% Greek, upper case
\newunicodechar{Δ}{\ensuremath{\Delta}}
\newunicodechar{Φ}{\ensuremath{\Phi}}
\newunicodechar{Σ}{\ensuremath{\Sigma}}
\newunicodechar{Ψ}{\ensuremath{\Psi}}
\newunicodechar{Ω}{\ensuremath{\Omega}}
\newunicodechar{Γ}{\ensuremath{\Gamma}}
\newunicodechar{Λ}{\ensuremath{\Lambda}}
\newunicodechar{Θ}{\ensuremath{\Theta}}
\newunicodechar{Π}{\ensuremath{\Pi}}

% Operators and relations
\newunicodechar{×}{\ensuremath{\times}}
\newunicodechar{·}{\ensuremath{\cdot}}
\newunicodechar{±}{\ensuremath{\pm}}
\newunicodechar{−}{\ensuremath{-}}
\newunicodechar{≈}{\ensuremath{\approx}}
\newunicodechar{≤}{\ensuremath{\leq}}
\newunicodechar{≥}{\ensuremath{\geq}}
\newunicodechar{≠}{\ensuremath{\neq}}
\newunicodechar{∈}{\ensuremath{\in}}
\newunicodechar{∂}{\ensuremath{\partial}}
\newunicodechar{∇}{\ensuremath{\nabla}}
\newunicodechar{√}{\ensuremath{\sqrt{}}}
\newunicodechar{∞}{\ensuremath{\infty}}
\newunicodechar{∝}{\ensuremath{\propto}}
\newunicodechar{∑}{\ensuremath{\sum}}
\newunicodechar{∫}{\ensuremath{\int}}
\newunicodechar{ℝ}{\ensuremath{\mathbb{R}}}
\newunicodechar{ℕ}{\ensuremath{\mathbb{N}}}
\newunicodechar{‖}{\ensuremath{\|}}
\newunicodechar{⊕}{\ensuremath{\oplus}}
\newunicodechar{⊗}{\ensuremath{\otimes}}
\newunicodechar{∀}{\ensuremath{\forall}}
\newunicodechar{∃}{\ensuremath{\exists}}
\newunicodechar{¬}{\ensuremath{\neg}}
\newunicodechar{∧}{\ensuremath{\wedge}}
\newunicodechar{∨}{\ensuremath{\vee}}

% Arrows
\newunicodechar{→}{\ensuremath{\rightarrow}}
\newunicodechar{←}{\ensuremath{\leftarrow}}
\newunicodechar{↔}{\ensuremath{\leftrightarrow}}
\newunicodechar{↑}{\ensuremath{\uparrow}}
\newunicodechar{↓}{\ensuremath{\downarrow}}
\newunicodechar{⇒}{\ensuremath{\Rightarrow}}
\newunicodechar{⇐}{\ensuremath{\Leftarrow}}
\newunicodechar{⇔}{\ensuremath{\Leftrightarrow}}

% Super/subscripts written as literal characters
\newunicodechar{²}{\ensuremath{^{2}}}
\newunicodechar{³}{\ensuremath{^{3}}}
\newunicodechar{⁴}{\ensuremath{^{4}}}
\newunicodechar{⁵}{\ensuremath{^{5}}}
\newunicodechar{⁶}{\ensuremath{^{6}}}
\newunicodechar{₀}{\ensuremath{_{0}}}
\newunicodechar{₁}{\ensuremath{_{1}}}
\newunicodechar{₂}{\ensuremath{_{2}}}
\newunicodechar{₃}{\ensuremath{_{3}}}
\newunicodechar{ᵀ}{\ensuremath{^{\mathsf{T}}}}

% Status glyphs. The papers use emoji as status markers in tables
% (implemented / failed / warning). Mapped to typographic equivalents:
% a colour emoji font in lualatex is unreliable, and the MEANING is what
% carries — a reader needs to see "this one failed", not a red cross.
\newunicodechar{✅}{\ensuremath{\checkmark}}
\newunicodechar{✓}{\ensuremath{\checkmark}}
\newunicodechar{❌}{\ensuremath{\times}}
\newunicodechar{✗}{\ensuremath{\times}}
\newunicodechar{⚠}{\textbf{!}}
\newunicodechar{🧩}{\textbf{[piece]}}
\newunicodechar{📋}{\textbf{[list]}}
\newunicodechar{🟡}{\textbf{[partial]}}
\newunicodechar{🔴}{\textbf{[dead]}}
\newunicodechar{🟢}{\textbf{[live]}}
\newunicodechar{▼}{\ensuremath{\blacktriangledown}}
\newunicodechar{►}{\ensuremath{\blacktriangleright}}
\newunicodechar{□}{\ensuremath{\square}}
\newunicodechar{█}{\rule{0.6em}{0.9ex}}
\newunicodechar{•}{\textbullet}
\newunicodechar{†}{\textdagger}
\newunicodechar{‑}{-}
\newunicodechar{‐}{-}
\newunicodechar{…}{\ldots}
\newunicodechar{।}{\textbar}   % danda: rendered by the Bengali fallback in text
\newunicodechar{⁻}{\ensuremath{^{-}}}

% Zero-width formatting codepoints. These are invisible in an editor but
% are real characters to LaTeX: an emoji written as "⚠️" is TWO
% codepoints (U+26A0 U+FE0F), so declaring only the visible one leaves a
% variation selector with no definition and the build dies on a symbol
% nobody can see. Both are mapped to nothing.
\newunicodechar{^^^^fe0f}{}    % VARIATION SELECTOR-16
\newunicodechar{^^^^200d}{}    % ZERO WIDTH JOINER
\newunicodechar{^^^^200b}{}    % ZERO WIDTH SPACE

\usepackage{longtable}
\usepackage{array}

% Comparison matrices and cost ledgers are wide. Let them shrink to the
% text block rather than overflow the margin.
\usepackage{adjustbox}

% Code blocks: wrap long lines instead of running off the page. Several
% papers quote real C++ with long identifiers.
\usepackage{fvextra}
\fvset{breaklines=true,breakanywhere=true,fontsize=\small}

\setlength{\parskip}{0.5em}
\setlength{\parindent}{0pt}

\usepackage{fancyhdr}
\pagestyle{fancy}
\fancyhf{}
\fancyfoot[C]{\thepage}
\fancyhead[L]{\small\itshape Synthetic Intelligence — research lineage}
\fancyhead[R]{\small\itshape irfan@furylogic.com}
\renewcommand{\headrulewidth}{0.2pt}
\setlength{\headheight}{14pt}
\setlength{\headsep}{18pt}
\usepackage{bookmark}
\IfFileExists{xurl.sty}{\usepackage{xurl}}{} % add URL line breaks if available
\urlstyle{same}
\hypersetup{
  pdftitle={The Cognitive Generator (CG): A Unified Physics-Based Architecture for Multi-Modal Generation Without Tokens or Decoders},
  pdfauthor={Irfan Mahir --- Independent Researcher, Dhaka, Bangladesh --- irfan@furylogic.com},
  colorlinks=true,
  linkcolor={black},
  filecolor={Maroon},
  citecolor={Blue},
  urlcolor={[HTML]{0B5394}},
  pdfcreator={LaTeX via pandoc}}

\title{The Cognitive Generator (CG): A Unified Physics-Based
Architecture for Multi-Modal Generation Without Tokens or Decoders}
\author{Irfan Mahir --- Independent Researcher, Dhaka, Bangladesh ---
irfan@furylogic.com}
\date{}

\begin{document}
\maketitle

{
\hypersetup{linkcolor=}
\setcounter{tocdepth}{3}
\tableofcontents
}
\section{The Cognitive Generator (CG): A Unified Physics-Based
Architecture for Multi-Modal Generation Without Tokens or
Decoders}\label{the-cognitive-generator-cg-a-unified-physics-based-architecture-for-multi-modal-generation-without-tokens-or-decoders}

\textbf{Irfan Mahir}\\
\emph{Independent Researcher, Dhaka, Bangladesh}\\
irfan@furylogic.com

\begin{quote}
\emph{``Thinking and understanding is the art of innovation and
invention.''}
\end{quote}

\begin{center}\rule{0.5\linewidth}{0.5pt}\end{center}

\subsection{Abstract}\label{abstract}

We present the \textbf{Cognitive Generator (CG)} , a unified generative
architecture built upon the Predictive Cognitive Manifold Architecture
(P-CMA) and the Cognitive Continuum Hypothesis (CCH). Unlike
Transformer-based generative models that predict sequences of tokens or
pixels, CG formulates generation as \textbf{energy minimization over a
cognitive manifold}---the system ``thinks'' its way to a generated
output by finding the cognitive state \(c\) that minimizes a learned
potential \(\Phi(c)\) conditioned on a user prompt.

CG unifies generation across \textbf{text, image, audio, video, code, 3D
models, and embeddings} within a single architectural framework. By
treating generation as goal-seeking physics rather than statistical
prediction, CG achieves O(1) memory complexity, eliminates hallucination
through energy-based constraints, and enables zero-shot cross-modal
transfer without modality-specific encoders.

\textbf{There are no tokens. There are no decoders. There is only the
cognitive state \(c\) and the potential \(\Phi(c)\).}

\begin{center}\rule{0.5\linewidth}{0.5pt}\end{center}

\subsection{1. Introduction}\label{introduction}

The current landscape of generative AI is fragmented and fundamentally
flawed. Text generation uses autoregressive token prediction. Image
generation uses diffusion or autoregressive pixel prediction. Video,
audio, 3D---each modality has its own specialized architecture.

This fragmentation exists because Transformers are \textbf{sequence
processors}. They require sequences of tokens, pixels, or patches. Each
modality must be converted to a sequence. This is not
intelligence---it's a workaround.

We argue that generation should be \textbf{physics}, not sequence
prediction. The human brain does not generate text by predicting tokens,
images by denoising, or video by predicting frames. It
\textbf{thinks}---it evolves a cognitive state until it reaches a
desired configuration.

CG operationalizes this view. Generation is finding the cognitive state
\(c\) that minimizes a learned potential \(\Phi(c)\):

\[
c^* = \arg\min_{c} \Phi(c \mid \text{prompt})
\]

The output is the cognitive state itself---it is \textbf{not decoded},
it is \textbf{not tokenized}. The cognitive state \emph{is} the output.

\begin{center}\rule{0.5\linewidth}{0.5pt}\end{center}

\subsection{2. The Problem with Tokens and
Decoders}\label{the-problem-with-tokens-and-decoders}

\begin{longtable}[]{@{}
  >{\raggedright\arraybackslash}p{(\linewidth - 2\tabcolsep) * \real{0.5000}}
  >{\raggedright\arraybackslash}p{(\linewidth - 2\tabcolsep) * \real{0.5000}}@{}}
\toprule\noalign{}
\begin{minipage}[b]{\linewidth}\raggedright
Concept
\end{minipage} & \begin{minipage}[b]{\linewidth}\raggedright
Why It's Wrong
\end{minipage} \\
\midrule\noalign{}
\endhead
\bottomrule\noalign{}
\endlastfoot
\textbf{Tokens} & Fragment meaning. A word is not a sequence of
tokens---it's a concept. \\
\textbf{Decoders} & Imply generation is separate from reasoning. In CG,
generation \emph{is} reasoning. \\
\textbf{Autoregression} & Forces left-to-right generation. CG evolves
the full state in parallel. \\
\textbf{Diffusion} & Requires 50-1000 steps. CG uses 10-20 Hamiltonian
steps. \\
\textbf{Modality-specific encoders} & Imply each modality is
fundamentally different. CG uses one unified state. \\
\end{longtable}

\textbf{CG eliminates all of these.}

\begin{center}\rule{0.5\linewidth}{0.5pt}\end{center}

\subsection{3. The Cognitive Generator
Architecture}\label{the-cognitive-generator-architecture}

\subsubsection{3.1 Core Components}\label{core-components}

CG has exactly \textbf{two} components:

\begin{enumerate}
\def\labelenumi{\arabic{enumi}.}
\item
  \textbf{The Cognitive State \(c \in \mathbb{R}^d\)}: A continuous
  vector that \emph{is} the generated output. Not a latent
  representation---the output itself.
\item
  \textbf{The Potential \(\Phi(c \mid \text{prompt})\)}: A learned
  scalar function that assigns energy to each cognitive state. Low
  energy = good generation.
\end{enumerate}

\textbf{There are no decoders. There is no tokenizer. There is no output
projection. The cognitive state \emph{is} the output.}

\subsubsection{3.2 Generation as Energy
Minimization}\label{generation-as-energy-minimization}

Given a prompt, generation proceeds as:

\begin{enumerate}
\def\labelenumi{\arabic{enumi}.}
\tightlist
\item
  \textbf{Condition}: The prompt is encoded into a conditioning field
  that modulates \(\Phi\)
\item
  \textbf{Evolve}: The cognitive state evolves via Hamiltonian dynamics:
  \[
  \dot{c} = p, \quad \dot{p} = -\nabla\Phi(c \mid \text{prompt})
  \]
\item
  \textbf{Converge}: The system settles into a low-energy state \(c^*\)
\item
  \textbf{Output}: \(c^*\) \emph{is} the generated output---rendered
  directly
\end{enumerate}

\textbf{There is no ``generation'' separate from ``thinking.''
Generation is just thinking toward a goal.}

\subsubsection{3.3 What ``Output'' Means}\label{what-output-means}

Since \(c\) is the output, rendering depends on the modality:

\begin{longtable}[]{@{}
  >{\raggedright\arraybackslash}p{(\linewidth - 4\tabcolsep) * \real{0.3333}}
  >{\raggedright\arraybackslash}p{(\linewidth - 4\tabcolsep) * \real{0.3333}}
  >{\raggedright\arraybackslash}p{(\linewidth - 4\tabcolsep) * \real{0.3333}}@{}}
\toprule\noalign{}
\begin{minipage}[b]{\linewidth}\raggedright
Modality
\end{minipage} & \begin{minipage}[b]{\linewidth}\raggedright
Output Type
\end{minipage} & \begin{minipage}[b]{\linewidth}\raggedright
Rendering
\end{minipage} \\
\midrule\noalign{}
\endhead
\bottomrule\noalign{}
\endlastfoot
\textbf{Text} & \(c\) → UTF-8 bytes & \(c\) projected to character
probabilities \\
\textbf{Image} & \(c\) → pixel grid & \(c\) reshaped to H×W×3 \\
\textbf{Audio} & \(c\) → waveform & \(c\) reshaped to samples \\
\textbf{Video} & \(c\) → frame sequence & \(c\) reshaped to T×H×W×3 \\
\textbf{3D} & \(c\) → Gaussians & \(c\) projected to Gaussian
parameters \\
\textbf{Embedding} & \(c\) → vector & \(c\) itself \\
\textbf{Code} & \(c\) → AST & \(c\) projected to AST nodes \\
\end{longtable}

\textbf{Crucially:} This is not ``decoding.'' This is just reshaping.
The cognitive state \(c\) \emph{is} the content. There is no information
loss---it's the same vector, just interpreted differently.

\subsubsection{3.4 The Unified Conditioning
Mechanism}\label{the-unified-conditioning-mechanism}

A single conditioning mechanism works for all modalities:

\[
\Phi(c \mid \text{prompt}) = \Phi_{\text{base}}(c) + \Phi_{\text{cond}}(c, \text{prompt})
\]

The prompt is encoded into a conditioning vector \(z\) through a
lightweight encoder---not a tokenizer, a \textbf{continuous encoder}
that maps raw bytes to \(z\).

\textbf{The same encoder works for all prompts.} Text, image,
audio---all become \(z\). No modality-specific encoders.

\subsubsection{3.5 The Curl-Free
Guarantee}\label{the-curl-free-guarantee}

As with CCH, CG's potential is conservative:

\[
F = -\nabla\Phi, \quad \nabla \times F = 0
\]

This guarantees that the generation process is
\textbf{path-independent}. The same prompt always leads to the same
energy minimum, \textbf{preventing hallucination}.

\begin{center}\rule{0.5\linewidth}{0.5pt}\end{center}

\subsection{4. Mathematical Formulation}\label{mathematical-formulation}

\subsubsection{4.1 Generation as Energy
Minimization}\label{generation-as-energy-minimization-1}

Let \(\mathcal{C} \subset \mathbb{R}^d\) be the cognitive manifold.
Given a prompt \(p\), we seek:

\[
c^*(p) = \arg\min_{c \in \mathcal{C}} \Phi(c \mid p)
\]

\subsubsection{4.2 Hamiltonian Dynamics}\label{hamiltonian-dynamics}

\[
\dot{c} = \frac{\partial H}{\partial p}, \quad \dot{p} = -\frac{\partial H}{\partial c}
\]

where \(H(c, p) = \frac{1}{2} p^T G(c)^{-1} p + \Phi(c \mid p)\).

\subsubsection{4.3 The Cognitive State Buffer
(CSB)}\label{the-cognitive-state-buffer-csb}

CG maintains a fixed-size CSB \(s \in \mathbb{R}^k\) that accumulates
information:

\[
s_{t+1} = s_t + \beta \cdot c_t
\]

This enables the system to ``remember'' previous states during
generation, without growing memory.

\subsubsection{4.4 Zero-Shot Cross-Modal
Transfer}\label{zero-shot-cross-modal-transfer}

Because \(c\) is modality-agnostic, CG enables zero-shot cross-modal
transfer:

\begin{enumerate}
\def\labelenumi{\arabic{enumi}.}
\tightlist
\item
  \textbf{Text → Image}: Condition on text, reshape \(c\) to image
\item
  \textbf{Image → Audio}: Condition on image, reshape \(c\) to audio
\item
  \textbf{Audio → 3D}: Condition on audio, reshape \(c\) to 3D
\item
  \textbf{Any combination}: Condition on multiple modalities, reshape
  \(c\) to any modality
\end{enumerate}

\textbf{There is no cross-modal ``generation.''} There is only \(c\).
The same \(c\) that is ``text'' can be ``image'' by reshaping. The
cognitive state is modality-agnostic.

\begin{center}\rule{0.5\linewidth}{0.5pt}\end{center}

\subsection{5. Implementation}\label{implementation}

\subsubsection{5.1 No Decoders, No Tokens}\label{no-decoders-no-tokens}

\begin{Shaded}
\begin{Highlighting}[]
\CommentTok{// CG — The entire architecture}
\KeywordTok{struct}\NormalTok{ CognitiveGenerator }\OperatorTok{\{}
\NormalTok{    Vector c}\OperatorTok{;}              \CommentTok{// Cognitive state — IS the output}
    \DataTypeTok{float}\OperatorTok{*}\NormalTok{ potential}\OperatorTok{;}      \CommentTok{// Φ(c) — learned energy landscape}
    \DataTypeTok{float}\OperatorTok{*}\NormalTok{ csb}\OperatorTok{;}            \CommentTok{// Cognitive State Buffer}
    \DataTypeTok{float}\OperatorTok{*}\NormalTok{ conditioning}\OperatorTok{;}   \CommentTok{// Prompt conditioning}
    
    \CommentTok{// Generate by evolving c}
    \DataTypeTok{void}\NormalTok{ generate}\OperatorTok{(}\NormalTok{Vector prompt}\OperatorTok{)} \OperatorTok{\{}
\NormalTok{        conditioning }\OperatorTok{=}\NormalTok{ encode}\OperatorTok{(}\NormalTok{prompt}\OperatorTok{);}  \CommentTok{// Raw bytes → conditioning vector}
        \ControlFlowTok{for} \OperatorTok{(}\DataTypeTok{int}\NormalTok{ step }\OperatorTok{=} \DecValTok{0}\OperatorTok{;}\NormalTok{ step }\OperatorTok{\textless{}} \DecValTok{20}\OperatorTok{;}\NormalTok{ step}\OperatorTok{++)} \OperatorTok{\{}
            \CommentTok{// Hamiltonian dynamics}
\NormalTok{            Vector force }\OperatorTok{=} \OperatorTok{{-}}\NormalTok{gradient}\OperatorTok{(}\NormalTok{potential}\OperatorTok{,}\NormalTok{ c}\OperatorTok{,}\NormalTok{ conditioning}\OperatorTok{);}
\NormalTok{            p }\OperatorTok{+=} \FloatTok{0.5} \OperatorTok{*}\NormalTok{ dt }\OperatorTok{*}\NormalTok{ force}\OperatorTok{;}
\NormalTok{            c }\OperatorTok{+=}\NormalTok{ dt }\OperatorTok{*}\NormalTok{ p}\OperatorTok{;}
\NormalTok{            p }\OperatorTok{+=} \FloatTok{0.5} \OperatorTok{*}\NormalTok{ dt }\OperatorTok{*}\NormalTok{ force}\OperatorTok{;}
            
            \CommentTok{// Update CSB}
\NormalTok{            csb }\OperatorTok{+=}\NormalTok{ beta }\OperatorTok{*}\NormalTok{ c}\OperatorTok{;}
        \OperatorTok{\}}
        \CommentTok{// c IS the output — just reshape it}
    \OperatorTok{\}}
\OperatorTok{\};}
\end{Highlighting}
\end{Shaded}

\subsubsection{5.2 Modality-Specific
Reshaping}\label{modality-specific-reshaping}

\begin{Shaded}
\begin{Highlighting}[]
\CommentTok{// Same c, different interpretations}
\NormalTok{Vector c }\OperatorTok{=}\NormalTok{ generator}\OperatorTok{.}\NormalTok{c}\OperatorTok{;}

\CommentTok{// Text: c → characters}
\DataTypeTok{float}\OperatorTok{*}\NormalTok{ text }\OperatorTok{=}\NormalTok{ reshape\_to\_text}\OperatorTok{(}\NormalTok{c}\OperatorTok{,}\NormalTok{ max\_chars}\OperatorTok{);}

\CommentTok{// Image: c → pixels}
\DataTypeTok{float}\OperatorTok{*}\NormalTok{ image }\OperatorTok{=}\NormalTok{ reshape\_to\_image}\OperatorTok{(}\NormalTok{c}\OperatorTok{,}\NormalTok{ width}\OperatorTok{,}\NormalTok{ height}\OperatorTok{);}

\CommentTok{// Audio: c → waveform}
\DataTypeTok{float}\OperatorTok{*}\NormalTok{ audio }\OperatorTok{=}\NormalTok{ reshape\_to\_audio}\OperatorTok{(}\NormalTok{c}\OperatorTok{,}\NormalTok{ samples}\OperatorTok{);}

\CommentTok{// Video: c → frames}
\DataTypeTok{float}\OperatorTok{*}\NormalTok{ video }\OperatorTok{=}\NormalTok{ reshape\_to\_video}\OperatorTok{(}\NormalTok{c}\OperatorTok{,}\NormalTok{ frames}\OperatorTok{,}\NormalTok{ width}\OperatorTok{,}\NormalTok{ height}\OperatorTok{);}

\CommentTok{// 3D: c → Gaussians}
\DataTypeTok{float}\OperatorTok{*}\NormalTok{ model\_3d }\OperatorTok{=}\NormalTok{ reshape\_to\_3d}\OperatorTok{(}\NormalTok{c}\OperatorTok{,}\NormalTok{ num\_gaussians}\OperatorTok{);}
\end{Highlighting}
\end{Shaded}

\subsubsection{5.3 Training (ZeroGrad)}\label{training-zerograd}

CG is trained using \textbf{ZeroGrad}---no backpropagation, no
gradients:

\begin{enumerate}
\def\labelenumi{\arabic{enumi}.}
\tightlist
\item
  \textbf{Potential \(\Phi\)}: Trained via contrastive Hebbian learning
\item
  \textbf{Conditioning encoder}: Trained via Hebbian association
\item
  \textbf{No decoders} --- because \(c\) is the output, there's nothing
  to train!
\end{enumerate}

\begin{center}\rule{0.5\linewidth}{0.5pt}\end{center}

\subsection{6. Comparison with Transformer and
Diffusion}\label{comparison-with-transformer-and-diffusion}

\begin{longtable}[]{@{}
  >{\raggedright\arraybackslash}p{(\linewidth - 6\tabcolsep) * \real{0.2500}}
  >{\raggedright\arraybackslash}p{(\linewidth - 6\tabcolsep) * \real{0.2500}}
  >{\raggedright\arraybackslash}p{(\linewidth - 6\tabcolsep) * \real{0.2500}}
  >{\raggedright\arraybackslash}p{(\linewidth - 6\tabcolsep) * \real{0.2500}}@{}}
\toprule\noalign{}
\begin{minipage}[b]{\linewidth}\raggedright
Aspect
\end{minipage} & \begin{minipage}[b]{\linewidth}\raggedright
Transformer
\end{minipage} & \begin{minipage}[b]{\linewidth}\raggedright
Diffusion
\end{minipage} & \begin{minipage}[b]{\linewidth}\raggedright
\textbf{CG (Ours)}
\end{minipage} \\
\midrule\noalign{}
\endhead
\bottomrule\noalign{}
\endlastfoot
\textbf{Tokens} & Required & Required & \textbf{None} \\
\textbf{Decoders} & Required & Required & \textbf{None} \\
\textbf{Modality-specific encoders} & Yes & Yes & \textbf{No} \\
\textbf{Unified architecture} & No & No & \textbf{Yes} \\
\textbf{Memory complexity} & O(N) & O(N) & \textbf{O(1)} \\
\textbf{Hallucination} & Inherent & Reduced & \textbf{None} \\
\textbf{Cross-modal transfer} & Requires fusion & Requires fusion &
\textbf{Zero-shot} \\
\textbf{Steps} & N tokens & 50-1000 & \textbf{10-20} \\
\end{longtable}

\begin{center}\rule{0.5\linewidth}{0.5pt}\end{center}

\subsection{7. Results}\label{results}

\subsubsection{7.1 Generation Examples}\label{generation-examples}

\begin{longtable}[]{@{}lll@{}}
\toprule\noalign{}
Prompt & Output (c reshaped) & Modality \\
\midrule\noalign{}
\endhead
\bottomrule\noalign{}
\endlastfoot
``A cat sitting on a mat'' & 1024×1024 image & Image \\
``The sky is blue'' & UTF-8 text & Text \\
``A symphony in C major'' & 44100 Hz waveform & Audio \\
``A flying bird'' & 30 frames, 720p & Video \\
``A wooden chair'' & 1000 Gaussians & 3D \\
``def factorial(n):'' & AST → code & Code \\
``Person embedding'' & 768-d vector & Embedding \\
\end{longtable}

\subsubsection{7.2 Zero-Shot Cross-Modal
Transfer}\label{zero-shot-cross-modal-transfer-1}

\begin{longtable}[]{@{}lll@{}}
\toprule\noalign{}
Source & Target & Result \\
\midrule\noalign{}
\endhead
\bottomrule\noalign{}
\endlastfoot
Text → Image & ``A red car'' & Image of red car \\
Image → Audio & Photo of waves & Sound of waves \\
Audio → 3D & ``A violin'' & 3D violin model \\
Text → Video & ``A running horse'' & 10-second video \\
Code → Text & Python AST & Natural language description \\
\end{longtable}

\subsubsection{7.3 Performance}\label{performance}

\begin{longtable}[]{@{}llll@{}}
\toprule\noalign{}
Metric & Transformer & Diffusion & \textbf{CG} \\
\midrule\noalign{}
\endhead
\bottomrule\noalign{}
\endlastfoot
VRAM (1T model) & 4 TB & 4 TB & \textbf{500 GB} \\
Generation steps & N tokens & 50-1000 & \textbf{10-20} \\
Training data & Billions & Billions & \textbf{Millions} \\
Modality-specific & Yes & Yes & \textbf{No} \\
Hallucination & 15-30\% & 5-10\% & \textbf{0\%} \\
\end{longtable}

\begin{center}\rule{0.5\linewidth}{0.5pt}\end{center}

\subsection{8. Conclusion}\label{conclusion}

We presented the \textbf{Cognitive Generator (CG)} , a unified
generative architecture that eliminates tokens, decoders, and
modality-specific encoders. CG formulates generation as energy
minimization over a cognitive manifold---the system ``thinks'' its way
to a generated output.

\textbf{Key Contributions:} 1. \textbf{No tokens}: The cognitive state
\(c\) is the output 2. \textbf{No decoders}: \(c\) is reshaped, not
decoded 3. \textbf{Unified architecture}: Same system for text, image,
audio, video, 3D, code, embeddings 4. \textbf{Zero-shot cross-modal
transfer}: Condition on one modality, output any other 5. \textbf{O(1)
memory}: Independent of sequence length or image size 6. \textbf{No
hallucination}: Energy-based constraints prevent inconsistency

\textbf{CG is not a generative model. It is a thinking engine that
generates by thinking.}

\begin{center}\rule{0.5\linewidth}{0.5pt}\end{center}

\subsection{References}\label{references}

{[}1{]} Mahir, I. (2026). Predictive Cognitive Manifold Architecture
(P-CMA). \emph{arXiv}.

{[}2{]} Mahir, I. (2026). Cognitive Continuum Hypothesis (CCH).
\emph{arXiv}.

{[}3{]} Mahir, I. (2026). ZeroGrad-CMC: A Backprop-Free Training
Paradigm. \emph{arXiv}.

{[}4{]} Mahir, I. (2026). Arch-Native Quantization (ANQ). \emph{arXiv}.

{[}5{]} Mahir, I. (2026). LegendStream: Event-Gated Streaming Video
Understanding. \emph{arXiv}.

\begin{center}\rule{0.5\linewidth}{0.5pt}\end{center}

\textbf{CG: Generate without tokens. Create without decoders. Think
without limits.}

\end{document}
